\begin{table}[ht]
\centering
\caption{Composite linguistic indices per chronological period. Values represent aggregate feature frequencies (per 1,000 words) across multiple linguistic dimensions. Archaism/Innovation indices measure retention vs. replacement of morphological features. Conservation Ratio $>1$ indicates archaic retention dominates innovation. Advanced indices quantify syntactic complexity, phonological conservation, and discourse sophistication.}
\label{tab:period-detection}
\small
\begin{tabular}{lrrrrrr}
\toprule
\textbf{Period} & \textbf{Arch.} & \textbf{Innov.} & \textbf{Cons.} & \textbf{Synt.} & \textbf{Text.} & \textbf{Overall} \\
 & \textbf{Index} & \textbf{Index} & \textbf{Ratio} & \textbf{Cplx.} & \textbf{Soph.} & \textbf{Innov.} \\
\midrule
Samhita (Early-Middle Vedic) & 23.1 & 20.6 & 1.13 & 2.8 & 5.1 & 3.2 \\
Brahmana (Late Vedic) & 27.2 & 24.5 & 1.13 & 6.7 & 13.1 & 7.6 \\
Upanishad (Latest Vedic) & 28.3 & 27.4 & 1.05 & 8.6 & 14.7 & 8.7 \\
Classical Sanskrit & 28.4 & 21.6 & 1.33 & 2.9 & 9.1 & 4.8 \\
\bottomrule
\end{tabular}
\end{table}